\subsection{Two minimal rows at each object and a stronger two-object bound}

\begin{proposition}\label{C:objectwise-two-row-cups}
Let $C$ be a finite strongly connected category with a
rectangular-band minimal ideal.  If each object has at most two
minimal rows, every product of rational degree-two classes on
$BC$ is zero.  The same statement holds if each object has at
most two minimal columns.  No finite-homotopy hypothesis is needed.
\end{proposition}
\begin{proof}
Use the row and column augmentation modules $V,W$ of
Lemma~\ref{C:category-band-Tor}.  A balanced functional decomposes
objectwise as $\lambda_i:W_i\otimes V_i\to\mathbb Q$.
Since $\dim V_i\leq1$, choose a row coordinate $a_i$ with values
zero and one when the row set has two elements, and value zero
when it has one.  Write $\lambda_i=b_i\otimes a_i$, normalizing
$b_i$ to be zero at a distinguished column point.  At a one-row
object choose $b_i=0$.

For $f:i\to j$, let $\chi_f$ be its row augmentation scalar,
zero if either row module is zero.  The row action and balancing
give affine equations
\[
 a_j\sigma_f=\chi_f a_i+a_f,
 \qquad b_i\tau_f=\chi_f b_j+b_f.
\]
These also cover the zero-dimensional endpoint cases.  If the
source has two rows and the target has one, balancing says that
$b_i\tau_f$ is constant; if the source has one row, $b_i=0$
makes the second equation automatic.  Composition gives
\[
 \chi_{fg}=\chi_f\chi_g,\qquad
 a_{fg}=a_f+\chi_f a_g,\qquad
 b_{fg}=b_g+\chi_g b_f.
\]
The band-grid cycles and their balancing three-chains in the proof
of Lemma~\ref{C:category-affine-balanced} apply separately at
each object.  They show that $u(f,g)=b_fa_g$ represents the
prescribed functional, and hence that every rational degree-two
class has this form.  That comparison does not require the
sum of the dimensions of the $V_i$ to be one.

Put $A_i=1/2$ at each two-row object and $A_i=0$ at each
one-row object.  Replace the offsets by
\[
 a'_f=a_f+\chi_f A_{s(f)}-A_{t(f)}.
\]
If $\chi_f\ne0$, the row map is a bijection between two-point
sets, so it preserves their midpoints and $a'_f=0$.
Consequently $\chi_fa'_f=0$ for every arrow.  The new offsets
still satisfy the typed affine composition equation.
For the normalized one-cochain $v(f)=A_{s(f)}b_f$ one has
\[
 \delta v(f,g)=b_f\bigl(A_{t(g)}-\chi_g A_{s(g)}\bigr),
 \qquad u'-u=-\delta v,
 \qquad u'(f,g)=b_fa'_g.
\]
The normalized three-cochain
$w(f,g,h)=b_fa'_g b_ga'_h$ satisfies
\[
 \delta w(f,g,h,k)
 =u'(f,g)u'(h,k)
  +b_fa'_k\,\chi_gb_g\,\chi_ha'_h
 =u'(f,g)u'(h,k).
\]
All equations are on composable strings, with
$s(f)=t(g)$ making the midpoint indices agree.
Every square therefore vanishes, and polarization over
$\mathbb Q$ gives every mixed product.  Passing to the opposite
category proves the assertion for columns.
\end{proof}

\begin{proposition}\label{C:one-active-three-by-three}
Let $C$ be a finite category with two objects $A,B$ and a
rectangular-band minimal ideal.  Suppose its row counts are
$(3,1)$ and its column counts are $(3,1)$, with both triples at
$A$.  If $BC$ has finite CW homotopy type, every product of two
rational degree-two classes on $BC$ is zero.
\end{proposition}
\begin{proof}
Let $H$ be the finite transformation-pair image of
$\operatorname{End}_C(A)$ on its three row and three column
points.  It contains the complete constant band
$\{x_{ij}:0\leq i,j\leq2\}$.  Put $e=x_{00}$ and let
$D$ be the full subcategory of $\operatorname{Kar}(H)$ on
objects $1,e$.
There is an explicit functor $q:C\to D$.
Send $A$ to $1$, $B$ to $e$, endomorphisms of $A$ to their
transformation pairs, and every endomorphism of $B$ to $e$.
An arrow $A\to B$ has the unique singleton row map and selects
a column $j$ at $A$; send it to $x_{0j}\in eH$.
An arrow $B\to A$ selects a row $i$ at $A$ and has the unique
singleton column map; send it to $x_{i0}\in He$.
The action equations and $x_{ij}x_{kl}=x_{il}$ check every
composable pair.  In particular, a composite through $A$ from
$B$ to itself maps to $x_{00}=1_e$, so identities are preserved.

The object $e$ is a retract of $1$, so $BH\to BD$ is a
classifying-space equivalence by the retract-generator lemma.
Moreover $q$ induces an isomorphism on $H_2(-;\mathbb Q)$.
Indeed $V_B=W_B=0$.  In the natural balanced tensor presentation
only $W_A\otimes V_A$ remains, and its only nonzero balancing
relations come from $\operatorname{End}_C(A)$; every cross-arrow
relation has a zero factor.  The presentation is therefore the
same as the one for $H$.  The band-grid cycles identify this
isomorphism with the actual functor-induced map.

Apply the complete three-by-three action-hull classification and
its exact certificates from the proof of
Theorem~\ref{B:three-by-three-cups}.  The maximal hull has exactly
the same balanced relation space as $H$, so every degree-two
class under consideration pulls back from that hull through
its two-object Karoubi category.  The rank-one, zero-product,
and fourth-homology-zero cases of that classification immediately
give the required vanishing after pullback.

We spell out the remaining skew-hull case, because finite type
is required only for the original $BC$ here.
The integral affine cocycle of that thirty-element hull extends
to the cohomology of its two-object Karoubi category through the
classifying-space equivalence.  Pull it back along $q$ to an
integral class on $BC$.
If one of its eight forbidden actors were the image of an element
of $\operatorname{End}_C(A)$, the eventual cyclic group of that
element would surject onto the actor's eventual cyclic group.
The nonzero integral character computed in
\eqref{B:3x3-cyclic-character} remains nonzero under this
surjection.  It is the restriction of the just-constructed class
on $BC$.  On the other hand, the local finite group maps to $C$
by adjoining a fresh identity if necessary, and Miller's theorem
makes the induced map to a finite CW model of $BC$ null-homotopic.
This is a contradiction.

Thus the actual image avoids all eight forbidden actors.  As
proved and independently checked in
Theorem~\ref{B:three-by-three-cups}, it lies in one of the nine
safe submonoids, each with zero fourth rational homology.  The
functor $q$ then factors through the corresponding Karoubi
category, and the pulled-back square vanishes.  Polarization
finishes the proof.  No finite-homotopy assertion about
$B\operatorname{End}_C(A)$, $BH$, or any quotient is used.
\end{proof}

\begin{theorem}[The strongly connected two-object exclusion through twenty-two]
\label{C:twentytwo-strongly-connected}
A finite category with two objects, nonempty hom-sets in both
directions, and at most twenty-two arrows cannot have classifying
space homotopy equivalent to $\mathbb{CP}^2$.
\end{theorem}
\begin{proof}
If either object identity lies in the minimal ideal, the retract
deletion in Lemma~\ref{C:finite-category-Rees} gives an equivalent
one-object category of smaller order.  This contradicts the monoid
minimum twenty-three in Theorem~\ref{B:CP2-minimum23}.  Thus both
identities survive outside the minimal ideal $K$.
Write $R=\sum_i r_i$ and $S=\sum_i c_i$ for its total row and
column counts.  If either augmentation module has dimension zero,
the band--Tor comparison gives $H_2=0$.

Suppose first that the minimal group $G$ is nontrivial.  Nonzero
$H_2$ requires $R,S\geq3$, and the arrow bound forces
$|G|=2$, $R=S=3$, and $|K|=18$.
The contracted category algebra has dimension at most four and
contains the two nonzero orthogonal object idempotents; hence it
is nonlocal.  Lemma~\ref{B:nonlocal-four-Tor} and the rational
comparison of Lemma~\ref{C:category-group-Tor} make a nonzero
fourth homology group incompatible with eventual homology
vanishing.  This excludes the group case.

Now $K$ is a band.  The existing rank-one obstruction in
Proposition~\ref{C:category-rank-one-cups} requires $R,S\geq4$.
If $RS\geq20$, at most two arrows lie outside $K$; they must be
exactly the two identities.  The contracted algebra is then
$\mathbb Q^2$, and the band--Tor formula gives $H_4=0$.
It follows that $R=S=4$.
Proposition~\ref{C:objectwise-two-row-cups} excludes a distribution
$(2,2)$ on either side.  Thus the row and column distributions
are both $(3,1)$, up to exchanging objects.
If the triples belong to different objects, each tensor
$W_i\otimes V_i$ is zero and $H_2=0$.
If they belong to the same object,
Proposition~\ref{C:one-active-three-by-three} gives zero
degree-two products.  Both conclusions contradict
$\mathbb{CP}^2$, completing the proof.
\end{proof}
